\section{Metrics and Diagnostic Categories}

\subsection{Persistence-relative skill}
Let \(y_h\) denote the verifying observation at horizon \(h\), and \(\hat{y}_h\) the corresponding point forecast produced by a model. Let \(\hat{y}^{(p)}_h\) denote the persistence baseline forecast for the same horizon. For a chosen loss function \(L(\cdot,\cdot)\), the baseline-relative (persistence-relative) skill score at horizon \(h\) is
\begin{equation}
\mathrm{Skill}(h) = 1 - \frac{\mathbb{E}[L(y_h,\hat{y}_h)]}{\mathbb{E}[L(y_h,\hat{y}^{(p)}_h)]}.
\end{equation}
In this work, the empirical results use an RMSE-based loss, so \(\mathrm{Skill}(h)>0\) indicates a lower RMSE than persistence for the same horizon and evaluation sample \cite{hyndman2006}.

Persistence-relative skill is useful because it anchors interpretation to an operationally meaningful comparator rather than to a model-to-model ranking alone. It also supports horizon-wise reading: in multi-step settings, skill can vary substantially across lead times.

At the same time, skill is a summary of pointwise error behaviour under a chosen loss. It does not, by itself, describe the trajectory shape of the point forecasts, and it does not indicate whether an apparent improvement is associated with reduced dynamic variability.

\subsection{Variance retention}
We define the variance-retention ratio at horizon \(h\) as
\begin{equation}
\alpha(h) = \frac{\sigma^2_{\hat{y}_h}}{\sigma^2_{y_h}}
          = \frac{\operatorname{Var}(\hat{y}_h)}{\operatorname{Var}(y_h)},
\end{equation}
the ratio of the predicted to the observed \emph{variance} at horizon \(h\).
Both variances are population variances (\(\mathrm{ddof}=0\)) computed over the
matched finite forecast--observation pairs within each station--model--horizon
cell, so \(\alpha(h)\) is the fraction of observed variance retained by the
forecast (a value of \(0.5\) corresponds to half the observed variance
retained). This is a variance ratio, not a standard-deviation ratio; the two
differ by a square and are not interchangeable.

The ratio \(\alpha(h)\) compares forecast trajectory variability with observed variability at the same horizon. Values \(\alpha(h)\ll 1\) indicate variance shrinkage consistent with smoothed point-forecast trajectories, while values \(\alpha(h)\gg 1\) indicate forecast variance inflation.

Variance retention should not be interpreted as a complete measure of trajectory quality. A forecast may retain the observed variance and still be biased, phase-shifted, or poor at reproducing individual episodes. The diagnostic is therefore intended to detect one specific failure mode—variance shrinkage—rather than to certify full dynamic adequacy.

Variance retention is not a calibration measure and is not intended to assess probabilistic uncertainty. It is a bounded post-evaluation diagnostic for reading point forecasts when operational interpretation depends on day-to-day fluctuations as well as on average accuracy \cite{jolliffe2012}.

\subsection{Auxiliary \texorpdfstring{\(\mathrm{Skill}_{VP}\)}{Skill\_VP}}
To summarise the joint behaviour of baseline-relative accuracy and variance retention, we also report an auxiliary variance-penalised skill diagnostic
\begin{equation}
\mathrm{Skill}_{VP}(h) = \mathrm{Skill}(h)\cdot\alpha(h).
\end{equation}

The auxiliary diagnostic is intended to highlight cases where positive persistence-relative skill co-occurs with strong variance shrinkage. The product form is not claimed to be uniquely optimal; it is used here only as a transparent descriptive penalty that keeps the sign of \(\mathrm{Skill}(h)\) while reducing positive skill when forecast variance is strongly collapsed. It should not be interpreted as a replacement for standard accuracy metrics, uncertainty quantification, or statistical significance testing \cite{diebold1995}.

Accordingly, \(\mathrm{Skill}_{VP}(h)\) is used only as a bounded diagnostic layer for post-evaluation interpretation. Its role is to provide a compact joint summary alongside \(\mathrm{Skill}(h)\) and \(\alpha(h)\), rather than to define a universal performance criterion.

\subsection{Diagnostic flags}
In addition to \(\alpha(h)\) and \(\mathrm{Skill}_{VP}(h)\), the post-evaluation reports diagnose flags for horizon-wise interpretation. The \texttt{collapse\_flag} is used when \(\alpha(h) < 0.5\). The \texttt{inflation\_flag} is used when \(\alpha(h) > 1.5\). The \texttt{near\_ideal\_flag} is used when \(\mathrm{Skill}(h)>0\) and \(0.8 \leq \alpha(h) \leq 1.2\). The \texttt{low\_sample\_flag} indicates fewer than 30 verification samples for a model/horizon cell.

These thresholds are diagnostic conventions for this post-evaluation, not universal decision rules. A collapse flag does not invalidate a model, and an inflation flag does not imply a better dynamic representation; the flags are intended to make horizon-wise interpretation auditable and reproducible.

